\documentclass[11pt]{article}
\usepackage[review]{../acl}
\usepackage{times}
\usepackage{latexsym}
\usepackage[T1]{fontenc}
\usepackage[utf8]{inputenc}
\usepackage{microtype}
\usepackage{amsmath,amssymb}
\usepackage{booktabs}
\usepackage{xurl}

\title{Mirror Cost: A Controlled Test of Reversed English}
\author{Anonymous submission}

\newcommand{\norm}[1]{n(#1)}
\newcommand{\rev}{\operatorname{rev}}
\newcommand{\seg}{\operatorname{seg}}

\begin{document}
% Full-paper snapshot 4; critique and successor changes are in FULL-PAPER-REVISION-LEDGER.md.
\maketitle

\begin{abstract}
We introduce mirror cost: both the forward letter stream and its reversal are segmented with the same vocabulary and procedure before causal-LM scoring. This matched treatment removes the direct spacing confound and measures the reverse-reading penalty.
\end{abstract}

\section{Introduction}

Constrained generation is usually described at the token level: a system must
include named terms, follow a schema, or satisfy a formal condition while
producing text \citep{hokamp2017grid,post2018fast,lu2022astar}. An exact
character-level palindrome is a harsher constraint. After case, spaces, and
punctuation are removed, a text $x$ must satisfy
\begin{equation}
  \norm{x}=\rev(\norm{x}).
\end{equation}
The first half of its letters determines the second half in reverse order. A
word boundary legal in one reading need not be legal in the other. The task is
therefore not merely to recognize a symmetric string: it is to ask whether one
letter sequence can be segmented into plausible English twice.

That distinction is visible in familiar examples. ``Madam, I'm Adam'' and
``Mad am I, madam'' have the same normalized letters but different word
boundaries and word orders. Long dictionary palindromes make the tension more
pronounced. Norvig's 90,439-letter version-3 result is an impressive exact
construction, but its author explicitly distinguishes length from plot and
character development \citep{norvig2016palindrome}. Length alone is not an
NLP quality metric.

We study the language cost before generating anything. Starting from ordinary
English, we remove formatting, segment the resulting letter stream in both
directions with the same vocabulary and objective, and score both readings
with the same causal language model. This symmetric control matters: comparing
the natural spacing of English with an automatically segmented reversal would
mix the cost of reversal with the cost of replacing the writer's boundaries.

Our contributions are:
\begin{itemize}
\item a bidirectional re-segmentation protocol for measuring the cost that an
  exact reversal constraint imposes on English;
\item a 57-cell empirical study showing a substantial reverse-reading
  likelihood penalty and loss of dictionary-word coverage; and
\item an exact two-ended decoder case study that separates structural candidate
  supply from linguistic quality. Its POS-shape gate is sound for its stated
  finite pattern language, but is not a grammar or a meaning model.
\end{itemize}
The paper makes no claim that the present system generates coherent long prose.

\section{Measuring the mirror cost}
\label{sec:measure}

Let $z=\norm{x}$ be the ASCII letter sequence of an English span and let $V$
be a fixed vocabulary. A segmentation procedure $s$ maps $z$ to a sequence of
vocabulary words, with single-letter fallbacks so that every stream can be
segmented. We use the same $V$ and $s$ for $z$ and $\rev(z)$. With a causal
language model $M$, define its per-letter cost as
\begin{equation}
  H_M(w_1\ldots w_k) =
  -\frac{\log_2 p_M(w_1\ldots w_k)}{|z|},
\end{equation}
where the denominator is the number of letters in the underlying stream, not
the number of model tokens. The \emph{mirror cost} is
\begin{equation}
  C_{M,s}(z)=H_M(\seg_s(\rev(z))) - H_M(\seg_s(z)).
\end{equation}
A positive value says that the same letters are less probable after reversal
and re-segmentation. It does not count grammatical strings, estimate a human
preference, or prove an information-theoretic lower bound.

\subsection{Data, vocabulary, and controls}

We draw spans of at least 20, 30, 40, 60, 80, or 120 letters at word boundaries
from the cached WikiText-2 training split \citep{merity2016pointer}. Each
configuration scores 150 spans. The vocabulary contains 52,927 lowercase
English headwords from the intersection of a dictionary and a frequency list.
Words have at most 20 letters; an unmatched letter is a legal one-character
fallback. We separately report coverage: the fraction of letters that land in
vocabulary words of length at least three. Thus a reversed stream cannot
appear well covered merely because a segmenter tiled it with one-letter units.

We test three segmentation objectives. \textsc{Unigram} maximizes summed Zipf
frequency minus a fixed per-word cost. \textsc{Fewest} minimizes the number of
segments. \textsc{Greedy} takes the longest available word from left to right,
with backtracking. These are deliberately different boundary preferences, not
claims that any is the best linguistic segmenter.

The retained results have 57 cells: DistilGPT-2, GPT-2, and GPT-2-medium
\citep{radford2019language} each score all three segmenters at six lengths;
GPT-2-large supplies an additional unigram check at 20, 30, and 40 letters.
All models are members of the GPT-2 family. They test scale sensitivity, not
independence across model architectures.

\section{Results: English is costly in reverse}
\label{sec:results}

Table~\ref{tab:unigram} gives the unigram-segmentation results, the condition
closest to a frequency-sensitive dictionary reading. The cost remains large as
spans grow. At 80 letters, the three complete model variants place 91\% of
forward letters but only 52\% of reversed letters inside dictionary words, and
pay 2.88--3.16 extra bits per letter for the reverse reading.

\begin{table}[t]
\centering
\small
\setlength{\tabcolsep}{3.2pt}
\begin{tabular}{rccrrr}
\toprule
Letters & Forward & Reverse & Distil & GPT-2 & Medium \\
 & coverage & coverage & cost & cost & cost \\
\midrule
20  & .90 & .48 & 3.38 & 3.44 & 3.60 \\
30  & .90 & .48 & 3.20 & 3.27 & 3.47 \\
40  & .90 & .50 & 3.03 & 3.16 & 3.30 \\
60  & .90 & .51 & 2.90 & 3.04 & 3.19 \\
80  & .91 & .52 & 2.88 & 3.03 & 3.16 \\
120 & .90 & .51 & 2.79 & 2.96 & 3.09 \\
\bottomrule
\end{tabular}
\caption{Unigram segmentation. Costs are additional bits per letter for the
reversed reading; each entry uses 150 WikiText-2 spans. GPT-2-large gives
3.60, 3.44, and 3.29 at 20, 30, and 40 letters, respectively.}
\label{tab:unigram}
\end{table}

The conclusion does not depend on the unigram objective alone. Across all 57
cells, mirror cost ranges from 2.16 to 3.60 bits per letter (mean 2.94).
The range by segmenter is 2.79--3.60 for \textsc{Unigram} (21 cells),
2.69--3.56 for \textsc{Fewest} (18), and 2.16--2.83 for \textsc{Greedy} (18).
The greedy procedure reduces the estimated gap, but does not remove it. Its
lower score is not evidence that a greedy segmentation makes reversed English
more grammatical; it only reveals that boundary choice affects the size of the
model-based diagnostic.

The coverage result is a model-free companion to the likelihood result. It
shows that roughly half the letters in the reversed spans cannot be placed in a
dictionary word of length three or more under these procedures. Coverage is
also not readability: a stream can be fully dictionary segmentable and still
be nonsensical.

\section{Exact decoding as a case study}
\label{sec:decoder}

The measurement explains why unconstrained left-to-right generation cannot
guarantee a palindrome: committing an early letter fixes a distant late letter
and its possible word boundaries. We use a two-ended overhang search derived
from the remainder construction of \citet{norvig2016palindrome}. A state
contains left and right word sequences plus the unmatched letters that one side
owes the other. A trie returns words compatible with that overhang. Every
closed candidate is checked by the normalized-letter predicate above, so model
scores can rank candidates but cannot make an invalid output valid.

The search can add structural constraints. Our finite POS-shape language is
derived from Brown-corpus universal tags \citep{francis1979brown,bird2009nltk,
petrov2012universal}. A partial left sequence must match the suffix of at
least one allowed tag pattern; a partial right sequence must match the prefix
of at least one allowed pattern. This is sound for the complete predicate:
future left additions only prepend to an already fixed suffix, and future right
additions only append to an already fixed prefix. A failed projection cannot
be repaired by a descendant. The gate is a finite-label viability test, not a
claim that an admitted word sequence is grammatical.

We ran 50 paired opening subtrees under a fixed 50,000-popped-state budget.
Opening selection and child order are independent of POS tags and outcomes;
both arms share vocabulary, letter bounds, terminal checks, and traversal
order. The terminal-only arm returns 407 accepted pairs; the incremental arm
returns 2,137. The incremental/terminal ratio is 5.47 for accepted pairs per
popped state (paired bootstrap 95\% interval 2.89--17.65) and 4.36 per CPU
second (2.30--13.93). It is 0.95 per generated state (0.46--3.20): pruning can
reach deeper, higher-branching regions within a pop budget and therefore
generate more candidate children. This is evidence about structural candidate
supply, not about the quality of the resulting text.

\section{Related work and scope}

\citet{papadopoulos2015palindromes} generate palindromes from n-gram corpora
with a graph that conjoins forward and backward constraints; their method also
uses syntactic patterns. We do not claim that prefix or suffix feasibility is
new in isolation, nor do we compare the two implementations. Our contribution
is instead the symmetric language measurement and its use to delimit what an
exact decoder result can show.

Lexically constrained decoding tracks unmet requirements during left-to-right
generation \citep{hokamp2017grid,post2018fast}; lookahead methods use future
feasibility to guide constrained text generation \citep{lu2022astar}. The
palindrome setting couples the two ends of the character string and allows word
boundaries to change under reversal. This makes symmetric re-segmentation a
useful control rather than a cosmetic preprocessing choice.

Our longer artifact illustrates the distinction between scale and language
quality. An audited inventory-aware adaptation reaches 90,937 letters, 498
more than Norvig's published version 3, while enforcing stated phrase and
repetition restrictions. It remains a dictionary list and is not a
matched-runtime comparison or a benchmark of meaning. We treat it as a stress
test for exact decoding, not as an NLP result about prose.

\section{Discussion and conclusion}

The empirical comparison is now controlled, but a single segmentation routine could manufacture its magnitude.

\section*{Limitations}

Mirror cost depends on a finite lexicon, a chosen segmenter, and a language model. It is not a grammar, semantic, or preference score.

\bibliography{../refs}

\appendix
\section{Reproduction}

The mirror-cost implementation is
\url{experiments/mirror_cost.py}; retained per-cell results are
\url{experiments/mirror_cost.json}. The program documents the three
segmentation objectives, corpus loading, fallback policy, and score
normalization. The decoder and its validator are under
\url{llm_palindrome/}; raw fixed-work trials and their audit are under
\url{runs/controlled-pos-pruning-openings-2026-09-11/}. The long artifact, its
phrase list, and an independent audit are under \url{artifacts/norvig-v3/}.
The released sources and results are enough to inspect the claims reported
here, but the missing sampled-span snapshot and heterogeneous model sweep make
the mirror-cost result a careful empirical study rather than a completed
benchmark.

\end{document}
